% Unidad U010. Espejo orientado y selección 3->1; redacción autónoma.

\chapter{Involución especular y selección local \(3\to1\) de la extensión superviviente}
\label{chap:espejo-seleccion-tres-uno}

La fase novena contiene dos estados locales de entrada. Una vez fijado el
estado orientado \(B_9\), la relación de frontera produce tres extensiones
hacia la fase primera. Las tres comparten un eje de autoescala y un agregado
ternario; difieren en el reparto ordenado de ese agregado entre las ramas de
clausura y propagación. Un segundo horizonte conserva exactamente una.

\section{Coordenadas simétrica y axial}

Sea
\[
 \mathcal B:=(\mathbb F_3^6)^3.
\]
Un bloque visible es
\[
 B=(u^\pi,u^e,u^\varphi)\in\mathcal B.
\]
Definimos
\[
 q(B):=u^\pi+u^e+u^\varphi,
 \qquad
 a(B):=u^\pi+u^e-u^\varphi.
\]
Todas las operaciones se realizan componente a componente en
\(\mathbb F_3^6\).

\begin{theorem}[Reconstrucción del eje y del agregado]
\label{thm:espejo-reconstruccion-eje-agregado}
Las coordenadas \(q,a\) determinan
\[
 \boxed{u^\varphi=2(q-a),}
 \qquad
 \boxed{u^\pi+u^e=q-u^\varphi.}
\]
\end{theorem}

\begin{proof}
Restar las definiciones da
\(q-a=2u^\varphi\). En \(\mathbb F_3\), \(2^{-1}=2\); por tanto
\(u^\varphi=2(q-a)\). Sustituir en la definición de \(q\) recupera el
agregado restante.
\end{proof}

La involución especular es
\[
 \mathfrak c(u^\pi,u^e,u^\varphi)
 =(u^e,u^\pi,u^\varphi).
\]

\begin{proposition}[Invariantes de la involución]
\label{prop:espejo-invariantes}
\(\mathfrak c^2=\operatorname{id}\), y
\[
 q(\mathfrak cB)=q(B),
 \qquad
 a(\mathfrak cB)=a(B).
\]
En particular, \(u^\varphi\) y \(u^\pi+u^e\) son invariantes.
\end{proposition}

\begin{proof}
Intercambiar dos veces las primeras coordenadas recupera \(B\). Tanto
\(q\) como \(a\) dependen de ellas sólo mediante su suma, de modo que el
intercambio las preserva. El \cref{thm:espejo-reconstruccion-eje-agregado}
transfiere esa invariancia al eje y al agregado.
\end{proof}

La involución no afirma que \(\pi=e\) ni que una sea el cociente de la otra.
Afirma que las dos ramas ocupan posiciones conjugadas respecto de una
estructura fija.

\section{Estado orientado en la fase \(9\) y datos de frontera asociados}

La proyección visible del estado seleccionado es
\[
 B_9=(100100\mid020112\mid010122).
\]
La relación inducida sobre bloques visibles se denota
\[
 \mathcal M_9^B\subseteq\mathcal B\times\mathcal B.
\]
Los dos estados locales de entrada de la tabla de fase preceden a esta
relación. Tras seleccionar el estado completo que proyecta a \(B_9\), su
fibra de salida tiene cardinal tres.

\begin{theorem}[Las tres extensiones de la fase novena]
\label{thm:espejo-tres-extensiones}
La imagen relacional de \(B_9\) contiene exactamente
\begin{align*}
B_{10}^{(0)}&=(101222\mid112221\mid112002),\\
B_{10}^{(1)}&=(102221\mid111222\mid112002),\\
B_{10}^{(2)}&=(111221\mid102222\mid112002).
\end{align*}
Las tres satisfacen
\[
 q_{10}=022112,
 \qquad
 a_{10}=101111,
\]
y, en consecuencia,
\[
 u_{10}^\varphi=112002,
 \qquad
 u_{10}^\pi+u_{10}^e=210110.
\]
\end{theorem}

\begin{proof}
La relación de extensión sobre estados completos se enumera en la fibra de
\(B_9\) y produce las tres filas indicadas. Sumar y restar sus tres
componentes da el mismo par \((q_{10},a_{10})\) en los tres casos. El
\cref{thm:espejo-reconstruccion-eje-agregado} produce el eje y el agregado
comunes. Las primeras dos componentes son distintas, por lo que las tres
filas representan distribuciones diferentes del mismo agregado.
\end{proof}

Sus publicaciones enteras en base tres son
\[
 (279,352,365),
 \qquad
 (280,351,365),
 \qquad
 (282,350,365).
\]

\section{Compatibilidad exacta de cilindros}

Una palabra ternaria de longitud \(L\) e índice \(N\) determina el cilindro
semiabierto
\[
 I_3(N,L)=\left[\frac{N}{3^L},\frac{N+1}{3^L}\right).
\]
Un prefijo de \(b\) bloques decimales, de entero asociado \(D\), determina
\[
 I_{1000}(D,b)=
 \left[\frac{D}{1000^b},\frac{D+1}{1000^b}\right).
\]

\begin{lemma}[Test entero de intersección]
\label{lem:espejo-test-entero-cilindros}
Los cilindros \(I_3(N,L)\) e \(I_{1000}(D,b)\) se intersectan si y sólo si
\[
 N1000^b<(D+1)3^L,
 \qquad
 (N+1)1000^b>D3^L.
\]
\end{lemma}

\begin{proof}
Dos intervalos semiabiertos \([a,b)\) y \([c,d)\) se intersectan si y sólo
si \(a<d\) y \(c<b\). Sustituimos los cuatro extremos y multiplicamos por
el entero positivo \(3^L1000^b\).
\end{proof}

El test utiliza sólo aritmética entera. No depende de una tolerancia de punto
flotante.

\section{2.º horizonte y unicidad}

El bloque siguiente es
\[
 B_{11}=(022212\mid110001\mid100021),
\]
y la frontera decimal correspondiente es
\[
 (502,662,638).
\]

\begin{theorem}[Resolución local \(3\to1\)]
\label{thm:espejo-seleccion-tres-uno}
Las tres extensiones del \cref{thm:espejo-tres-extensiones} son compatibles
en profundidad uno. Al exigir simultáneamente \(B_{11}\) y la frontera
\((502,662,638)\), sólo \(B_{10}^{(0)}\) se prolonga como el mismo estado
completo. Por tanto,
\[
 \boxed{
 \#\operatorname{Surv}^{\rm cal}_1(B_9)=3,
 \qquad
 \#\operatorname{Surv}^{\rm cal}_2(B_9)=1.}
\]
\end{theorem}

\begin{proof}
Aplicamos el \cref{lem:espejo-test-entero-cilindros} a cada uno de los tres
canales. En el primer horizonte, las seis desigualdades de cada candidato se
satisfacen. Tras anexar \(B_{11}\), las desigualdades de los tres canales
se mantienen para \(B_{10}^{(0)}\). En \(B_{10}^{(1)}\) y
\(B_{10}^{(2)}\) se mantiene la compatibilidad del canal \(\varphi\), pero
fallan las condiciones de las ramas \(\pi\) y \(e\). La enumeración deja,
por tanto, un único estado a profundidad dos.
\end{proof}

Esta prueba no define la conexión global a partir de un ejemplo local. Es un
testigo explícito de cómo una relación multivaluada en un horizonte puede
volverse univaluada al conservar memoria y exigir una prolongación mayor.

\section{Retorno de fase con cambio de bloque}

La fase primera de la vuelta inicial es
\[
 B_1=(012222\mid121221\mid112202),
\]
mientras la fase primera de la vuelta siguiente es
\[
 B_{10}=(101222\mid112221\mid112002).
\]
Así,
\[
 g(1)=g(10)=1,
 \qquad
 B_1\neq B_{10}.
\]
La fase retorna, pero el eje, el agregado, el reparto y la memoria de
cociente han evolucionado.

\section{Diagrama de la involución especular \(\pi/e\), el eje fijo \(\varphi\) y la selección local \(3\to1\) por prolongabilidad}

\begin{figure}[htbp]
\centering
\resizebox{0.98\textwidth}{!}{%
\begin{tikzpicture}[
 node distance=10mm and 15mm,
 every node/.style={font=\small,align=center},
 state/.style={draw=black,rounded corners,inner sep=5pt},
 inv/.style={draw=black,double,rounded corners,inner sep=5pt},
 arr/.style={-{Latex[length=2mm]},semithick}
]
 \node[state] (pi) {$u^\pi$\\rama de clausura};
 \node[state,right=70mm of pi] (e) {$u^e$\\rama de propagación};
 \coordinate (mirroraxis) at ($(pi)!0.5!(e)$);
 \node[inv,above=9mm of mirroraxis] (phi) {$u^\varphi$\\eje fijo};
 \draw[arr,<->] (pi) -- node[below] {$\mathfrak c$} (e);
 \node[above=8mm of phi] (aggregate)
   {$u^\pi+u^e$ conservado en la fibra};
 \draw[dashed] (phi.north) -- (aggregate.south);
 \node[state,below=20mm of mirroraxis] (b9) {fase $9$: $B_9$\\dos entradas locales};
 \node[state,below left=18mm and 25mm of b9] (b100) {$B_{10}^{(0)}$};
 \node[state,below=18mm of b9] (b101) {$B_{10}^{(1)}$};
 \node[state,below right=18mm and 25mm of b9] (b102) {$B_{10}^{(2)}$};
 \node[state,below=18mm of b101] (b11) {$B_{11}$\\única prolongación};
 \draw[arr] (b9) -- (b100);
 \draw[arr] (b9) -- (b101);
 \draw[arr] (b9) -- (b102);
 \draw[arr] (b100) -- (b11);
 \draw[arr,densely dotted] (b101) -- node[right] {obstruida} (b11);
\draw[arr,densely dotted] (b102) -- node[right] {obstruida} (b11);
\end{tikzpicture}
}
\caption{El espejo conserva el eje \(u^\varphi\) y el agregado
\(u^\pi+u^e\). El retorno de la fase novena abre tres extensiones visibles;
la memoria del segundo horizonte conserva una. La vuelta no reinicia el
estado.}
\label{fig:espejo-tres-uno-arrastre}
\end{figure}

\section{Memoria residuo--cociente durante el espejo}

La recurrencia
\[
 x_kA=u_k+3x_{k+1}
\]
de la \cref{eq:bockstein-hensel-local} se aplica a cada uno de los tres
canales. La parte \(u_k\) publica el bloque visible; \(x_{k+1}\) conserva
el cociente. Como \(A\in\operatorname{GL}_6(\mathbb Z)\), el par recupera
el estado anterior. El espejo puede intercambiar las ramas visibles sin
eliminar sus cocientes de procedencia.

\section{Obstrucciones de la involución especular \(\pi/e\) y de la selección local \(3\to1\): invariantes, horizonte de profundidad \(2\) y memoria Hensel}

\begin{criterion}[Criterios de refutación del espejo]
La construcción queda refutada si:
\begin{enumerate}
 \item la involución modifica \(u^\varphi\) o \(u^\pi+u^e\);
 \item se confunden los dos estados locales de entrada con las tres salidas;
 \item alguna de las tres salidas tiene distinto \((q_{10},a_{10})\);
 \item el test de cilindros utiliza una tolerancia flotante;
 \item más de una salida sobrevive al segundo horizonte declarado;
 \item el retorno de fase se interpreta como igualdad \(B_1=B_{10}\);
 \item el intercambio visible borra los cocientes de Hensel.
\end{enumerate}
\end{criterion}

El espejo y la selección local están ahora construidos. Las unidades
siguientes desarrollarán las vacancias sturmianas y, después, las tres
biografías analíticas completas.
